\documentclass[12pt]{article}
\usepackage[utf8]{inputenc}
\usepackage{amsmath, amssymb, amsthm}
\usepackage[margin=1in]{geometry}

\newtheorem{theorem}{Theorem}
\newtheorem{lemma}[theorem]{Lemma}
\newtheorem{definition}{Definition}

\title{Problem 300: Protein Folding}
\author{Project Euler Solution}
\date{}

\begin{document}
\maketitle

\section{Problem Statement}
A protein of length $n$ is a string over $\{H, P\}$. It folds as a self-avoiding walk (SAW) on the 2D square lattice. An \textbf{H-H contact}: two H residues at non-consecutive positions that are lattice-adjacent. Each protein folds optimally (maximizing contacts).

Reference: $n = 8$ gives average $850/256 = 3.3203125$. Find the average for $n = 15$ (exact decimal).

\section{Mathematical Foundation}

\begin{definition}
A \emph{self-avoiding walk} of length $n-1$ on $\mathbb{Z}^2$ is an injective sequence $(v_0, \ldots, v_{n-1})$ with $|v_{i+1} - v_i|_1 = 1$. Its \emph{contact topology} is $\mathcal{C} = \{(i,j) : |i-j| > 1,\, |v_i - v_j|_1 = 1\}$.
\end{definition}

\begin{theorem}[Decomposition of Contacts]
Total contacts $= \mathrm{consec}(p) + \mathrm{non\text{-}consec}(p, w)$, where $\mathrm{consec}(p) = |\{i : p_i = p_{i+1} = H\}|$ depends only on $p$, and $\mathrm{non\text{-}consec}$ depends on the walk.
\end{theorem}

\begin{proof}
Consecutive residues are always lattice-adjacent in any SAW, so their contacts depend only on the protein string. Non-consecutive contacts depend on the walk geometry via the contact topology.
\end{proof}

\begin{theorem}[Optimality via Topology Enumeration]
\[
\mathrm{opt}(p) = \mathrm{consec}(p) + \max_{\mathcal{C} \in \mathcal{T}} \sum_{(i,j) \in \mathcal{C}} \mathbf{1}[p_i = H \wedge p_j = H]
\]
where $\mathcal{T}$ is the set of all realizable contact topologies.
\end{theorem}

\begin{proof}
Since $\mathrm{consec}(p)$ is constant across walks, optimization reduces to maximizing non-consecutive contacts. These depend only on the contact topology, so we enumerate distinct topologies.
\end{proof}

\begin{lemma}[Bitmask Evaluation]
With protein bitmask $B$ (bit $i = 0$ iff $p_i = H$) and contact mask $M_{ij} = (1 \ll i) | (1 \ll j)$, the pair $(i,j)$ contributes iff $(B \;\&\; M_{ij}) = 0$.
\end{lemma}

\begin{proof}
Both $p_i = H$ and $p_j = H$ iff bits $i$ and $j$ of $B$ are both 0, i.e., $B \;\&\; M_{ij} = 0$.
\end{proof}

\begin{lemma}[Symmetry Reduction]
The dihedral group $D_4$ of the square lattice preserves contact topologies. Fixing the first step to go rightward reduces enumeration by a factor of 4.
\end{lemma}

\begin{proof}
Rotation by $90^\circ$ maps SAWs to SAWs with identical contact topologies (lattice adjacency is rotation-invariant). Fixing the first step direction eliminates 4 rotational equivalences.
\end{proof}

\begin{theorem}[Exact Rationality]
$E[\mathrm{opt}(p)] = \frac{1}{2^n} \sum_p \mathrm{opt}(p) \in \frac{1}{2^n}\mathbb{Z}$, hence the answer is an exact terminating decimal.
\end{theorem}

\begin{proof}
Each of the $2^n$ proteins is equally likely, and each $\mathrm{opt}(p)$ is a non-negative integer. Thus the expectation has denominator dividing $2^n = 2^{15} = 32768$, which is a power of 2, ensuring a terminating decimal.
\end{proof}

\section{Algorithm}

\begin{verbatim}
1. Enumerate SAWs of 14 steps (first step fixed rightward)
2. Extract contact topologies, deduplicate
3. For each protein bitmask B in {0, ..., 2^15 - 1}:
   a. Compute consec(B)
   b. For each topology C, count non-consec contacts via bitmask
   c. opt(B) = consec(B) + max over C
4. Return sum of opt(B) / 2^15
\end{verbatim}

\section{Complexity Analysis}
\begin{itemize}
\item \textbf{Time:} $O(W + |\mathcal{T}| \cdot 2^n)$ where $W \sim 10^6$ (SAWs after symmetry), $|\mathcal{T}| \sim 5000$ topologies, $2^{15} = 32768$ proteins. Total $\approx 1.6 \times 10^8$ operations.
\item \textbf{Space:} $O(|\mathcal{T}| \cdot \bar{c})$ for topologies ($\bar{c}$ = avg contacts per topology), plus $O(n)$ for backtracking.
\end{itemize}

\section{Answer}

\[
\boxed{8.0540771484375}
\]

\end{document}
